\documentclass[10pt,a4paper]{amsart}
\usepackage[margin=19mm]{geometry}
\usepackage{amsmath,amssymb,mathtools}
\usepackage[scr=rsfs,cal=euler]{mathalfa}
\usepackage{xfrac}
\usepackage[unicode,hidelinks,bookmarks=false]{hyperref}
\newcommand{\ie}{i.e.\ }
\newcommand{\cf}{cf.\ }
\newcommand{\resp}{respectively\ }
\newcommand{\Subsection}{\subsection}
\providecommand{\nobreakdash}{\nobreak\hbox{-}}
\setlength{\emergencystretch}{2em}
\multlinegap0pt
\usepackage{bm}
\usepackage{bbm}
\usepackage{dsfont}
\usepackage{xspace}
\usepackage{cancel}
\usepackage{mathtools}
\usepackage{amsthm}
\makeatletter
\@ifundefined{theorem}{\newtheorem{theorem}{Theorem}}{}
\@ifundefined{lemma}{\newtheorem{lemma}[theorem]{Lemma}}{}
\makeatother
\newcommand{\F}{\mathbb{F}}
\newcommand{\fl}{\mathfrak{l}}
\newcommand{\op}[1]{\operatorname{#1}}
\title[Galois representations are surjective for almost all Drinfel]{Galois representations are surjective for almost all Drinfeld modules}
\author{Anwesh Ray}
\date{Source: arXiv:2407.14264v2 (CC BY 4.0)}
\begin{document}
\maketitle

\noindent\emph{Source and licence.} Galois representations are surjective for almost all Drinfeld modules. Version arXiv:2407.14264v2 (CC BY 4.0), distributed under the Creative Commons Attribution 4.0 International licence, as recorded in the arXiv licence metadata for this exact version and shown on the abstract page https://arxiv.org/abs/2407.14264v2. Full rights notice: licenses/arxiv-2407.14264-CC-BY-4.0.txt.\par\smallskip

\noindent\textit{Editorial scope.} Counted unit: the Lemma whose source label is \texttt{v equals technical lemma} in the article (source0.tex lines 560--567, source label \texttt{v equals technical lemma}), delivered together with the article's own complete original proof of it (source0.tex lines 568--585, one \texttt{proof} environment of 18 lines). No mathematics is written, repaired or re-derived: statement and proof are the article's own text line for line. Required context retained uncounted: tiny lemma (lines 552--554). Every definition, display and label of those lines is retained unchanged. Omitted from this package and not redistributed: 1--551, 555--559, 586--660. Nothing inside a retained excerpt is omitted or altered: every retained body is delivered exactly as the article's lines are written, so no omission receipt of any kind accompanies this package. Citations: the retained excerpts carry no citation command at all, so no bibliography entry of the article is transcribed and no reference is redistributed. Reference adaptation: none was needed and none was made; every reference inside the delivered unit resolves inside it, so no unresolved reference or citation is produced. Printed slips kept literal: no printed word, symbol, hypothesis or number of the counted statement or of its proof was altered. Layout and class: the article's own class is \texttt{amsart} with packages including geometry, amscd, amsmath, amsxtra, amsthm, amssymb, stmaryrd, xr, mathrsfs, mathtools, enumerate, commath, comment, enumitem; the wrapper is the lane's pinned \texttt{amsart} editorial wrapper at 10pt on A4, which restates the theorem-like environments the retained text needs and adds the article's own macro definitions that the retained text needs (\texttt{\textbackslash F}, \texttt{\textbackslash fl}, \texttt{\textbackslash op}), with extra wrapper packages (bm, bbm, dsfont, xspace, cancel, mathtools). The delivered numbering is the unit's own. Assets: no retained span contains a figure environment, a graphics-inclusion command, a PDF-page inclusion, a table or a TikZ picture; no image file is copied into this package, the package declares no asset dependency and no image file is redistributed with it. Licence: version arXiv:2407.14264v2 of this article is distributed under the Creative Commons Attribution 4.0 International licence, as recorded in the arXiv licence metadata for this exact version and shown on the abstract page https://arxiv.org/abs/2407.14264v2. A full rights notice is delivered at licenses/arxiv-2407.14264-CC-BY-4.0.txt. Characters: every retained excerpt is pure ASCII, so the delivered engine, XeLaTeX with its default Latin Modern text font, reports no missing character.
\begin{lemma}\label{tiny lemma}
    Let $w\in F_{\fl}^{\op{sep}}$ be such that $v(w)\geq 0$. Then we find that $v(e_\Lambda(w))=v(w)$. 
\end{lemma}

\begin{lemma}\label{v equals technical lemma}
    With respect to notation above, we have that 
    \[\begin{split} & v\left(e_\Lambda(a_1 w_1+a_2 w_2+\dots+a_{r-1} w_{r-1} +b z)\right) \\ = &\begin{cases}
        q^{(r-1)i}v(z) & \text{ if }b\neq 0\text{ and }\op{deg}_T(b)=i.\\
        v(a_1 w_1+a_2 w_2+\dots+a_{r-1} w_{r-1})& \text{ otherwise. }
    \end{cases}
    \end{split}\]
\end{lemma}

\begin{proof}
    \par First we consider the case when $b\neq 0$ and write $b=c+dT$, where $c, d\in \F_q$. One finds that 
\[\begin{split} & e_\Lambda(a_1 w_1+a_2 w_2+\dots+a_{r-1} w_{r-1} +b z) \\ 
=& \sum_{i=1}^{r-1} \varphi_{a_i} \left(e_{\Lambda}(w_i)\right)+\varphi_b \left(e_{\Lambda}(z)\right).\end{split}\]
Note that $w_i$ is a root of $\varphi_{T^2}(x)$. We write \[\varphi_T(\tau)=T+b_1 \tau+ b_2 \tau^2+\dots+ b_{r-1}\tau^{r-1}\] and note that $v(b_{r-1})=0$ since $\varphi$ has good reduction. We find that 
\[\varphi_{T^2}=\varphi_T\circ \varphi_T=T^2+c_1\tau+\dots +c_{2(r-1)}\tau^{2(r-1)},\] where $c_{2(r-1)}=b_{r-1}^{1+q^{r-1}}$. In particular, $v(c_{2(r-1)})=0$ and all slopes of the Newton polygon of $\phi_{T^2}(x)$ are $\leq 0$. We deduce therefore that all solutions to $\varphi_{T^2}(x)=0$ have non-negative valuation, in particular, $v(w_i)\geq 0$ for $1\leq i\leq r-1$. In particular, we find that $v(w_i)\geq 0$. It follows from Lemma \ref{tiny lemma} that $v(e_{\Lambda}(w_i))=v(w_i)\geq 0$. We find that
\[\varphi_b(x)=cx+d\varphi_T(x)=bx+db_1x^q+\dots +d b_{(r-1)} x^{q^{(r-1)}},\]
and therefore, 
\[v\left(\varphi_b(e_\Lambda(z))\right)=q^{(r-1)}v(e_\Lambda(z))=\begin{cases}
    q^{(r-1)}v(z) &\text{ if }d\neq 0;\\
    v(z) &\text{ if }d=0.
\end{cases}\]
Therefore, we have shown that if $b\neq 0$, 
\[v\left(e_\Lambda(a_1 w_1+a_2 w_2+\dots+a_{r-1} w_{r-1} +b z)\right)=v\left(\varphi_b(e_\Lambda(z))\right)=q^{(r-1)i}v(z),\]
where $i:=\op{deg}_T(b)$.
\par When $b=0$, it follows from Lemma \ref{tiny lemma} that 
\[v\left(e_\Lambda(a_1 w_1+a_2 w_2+\dots+a_{r-1} w_{r-1})\right)=v(a_1 w_1+a_2 w_2+\dots+a_{r-1} w_{r-1}).\]
\end{proof}

\end{document}
